
%  GUÍA PRÁCTICA DE FUNCIONES — Usach Premium
%  Cálculo I | Progresión básico → PEP
%  Resultados del solucionario validados con SymPy + NumPy (36/36 PASS)
%
%  Compila con:
%    pdflatex -interaction=nonstopmode guia-practica-funciones.tex   (2 veces)
% ============================================================

\documentclass[letterpaper, 12pt]{article}

% ── Paquetes ─────────────────────────────────────────────────
\usepackage[utf8]{inputenc}
\usepackage[spanish]{babel}
\usepackage{amsmath, amssymb}
\usepackage{geometry}
\usepackage{fancyhdr}
\usepackage{enumitem}
\usepackage{graphicx}
\usepackage{hyperref}
\usepackage{multicol}
\usepackage{parskip}
\usepackage{xcolor}
\usepackage{tcolorbox}
\usepackage{eso-pic}
\usepackage{transparent}
\usepackage{colortbl}
\usepackage{booktabs}
\usepackage{tikz}
\usetikzlibrary{patterns, arrows.meta, babel}
\tcbuselibrary{skins, breakable}

% ── Geometría ────────────────────────────────────────────────
\geometry{letterpaper, margin=1.5cm, top=3cm, bottom=2.5cm, headheight=2cm}

% ── Colores institucionales ──────────────────────────────────
\definecolor{celesteSuave}{HTML}{F0F7FF}
\definecolor{azulFuerte}{HTML}{1A5276}
\definecolor{turquesaUsach}{HTML}{33CCCC}
\definecolor{enlace}{HTML}{1A5276}
\definecolor{verdePaso}{HTML}{1E8449}
\definecolor{rojoAlerta}{HTML}{C0392B}
\definecolor{naranjaNivel}{HTML}{D35400}

% ── Rutas de logos ───────────────────────────────────────────
\newcommand{\logoup}{./images/logo_up.png}
\newcommand{\logousach}{./images/logo_usach.png}
\newcommand{\logofondo}{./images/logo_up.png}

% ── Información de la guía ───────────────────────────────────
\newcommand{\institucion}{Usach Premium}
\newcommand{\curso}{Cálculo I}
\newcommand{\guiasnumero}{Guía Práctica — Funciones}
\newcommand{\semestre}{1er.\ Semestre de 2026}
\newcommand{\preparadopor}{Usach Premium.}
\newcommand{\webinstitucion}{www.usachpremium.cl}
\newcommand{\instagraminstitucion}{https://www.instagram.com/usach.premium}
\newcommand{\transparencia}{0.05}
\newcommand{\fraseportada}{``Por y para estudiantes''}

\newcommand{\listacontenidos}{
    \item[\color{azulFuerte}$\Rightarrow$] Dominio y recorrido de funciones
    \item[\color{azulFuerte}$\Rightarrow$] ¿Dónde crece y dónde decrece la función?
    \item[\color{azulFuerte}$\Rightarrow$] Composición de funciones
    \item[\color{azulFuerte}$\Rightarrow$] Inyectividad y sobreyectividad
    \item[\color{azulFuerte}$\Rightarrow$] Biyectividad y función inversa
    \item[\color{azulFuerte}$\Rightarrow$] Ejercicios tipo PEP
}

% ── Hyperref ─────────────────────────────────────────────────
\hypersetup{colorlinks=true,linkcolor=enlace,urlcolor=enlace,pdfborder={0 0 0}}

% ── Encabezado y pie ─────────────────────────────────────────
\pagestyle{fancy}
\fancyhf{}
\fancyhead[L]{\includegraphics[height=1.2cm]{\logoup}}
\fancyhead[R]{%
    \begin{minipage}[b]{0.45\textwidth}
        \raggedleft\fontsize{9pt}{11pt}\selectfont
        \textbf{\color{azulFuerte}\institucion} \\
        \curso\ — \guiasnumero \\
        \textit{\color{gray}@usach.premium}
    \end{minipage}\hspace{0.1cm}%
    \includegraphics[height=1.2cm]{\logousach}%
}
\fancyfoot[L]{\fontsize{9pt}{11pt}\selectfont \textit{\fraseportada}}
\fancyfoot[R]{\fontsize{9pt}{11pt}\selectfont Página \thepage}
\renewcommand{\headrulewidth}{0.4pt}
\renewcommand{\footrulewidth}{0.4pt}

\fancypagestyle{plain}{\fancyhf{}\renewcommand{\headrulewidth}{0pt}\renewcommand{\footrulewidth}{0pt}}

% ── Marca de agua ────────────────────────────────────────────
\newcommand{\MarcaAgua}{%
    \AddToShipoutPicture{%
        \AtPageCenter{\makebox(0,0){%
            \transparent{\transparencia}%
            \includegraphics[width=0.7\textwidth]{\logofondo}}}}}

% ── Cajas ────────────────────────────────────────────────────
\newenvironment{kitpropiedades}[1]{%
    \begin{tcolorbox}[colback=celesteSuave, colframe=azulFuerte,
        title=\textbf{#1}, fonttitle=\bfseries]\small
}{\end{tcolorbox}}

\newtcolorbox{desafiobox}[1]{%
    colback=white, colframe=azulFuerte, fonttitle=\bfseries,
    coltitle=white, title=#1, arc=8pt, boxrule=1.2pt, enhanced,
    drop shadow={black!5}}

\newenvironment{solbox}[1]{%
    \begin{tcolorbox}[colback=celesteSuave, colframe=azulFuerte,
        title=\textbf{#1}, fonttitle=\bfseries]
}{\end{tcolorbox}}

% Caja de nivel
\newcommand{\nivelbasico}{\textcolor{verdePaso}{\textbf{[Básico]}}}
\newcommand{\nivelintermedio}{\textcolor{naranjaNivel}{\textbf{[Intermedio]}}}
\newcommand{\nivelpep}{\textcolor{rojoAlerta}{\textbf{[PEP]}}}

\newenvironment{soluciones}{%
    \newpage
    \begin{center}
        \begin{tcolorbox}[colback=azulFuerte,colframe=azulFuerte,
            colupper=white,width=0.8\textwidth,center,arc=10pt]
            \centering\Large\textbf{SOLUCIONARIO}
        \end{tcolorbox}
    \end{center}
    \MarcaAgua\small
}{}

% ════════════════════════════════════════════════════════════
\begin{document}
\thispagestyle{plain}

% ── PORTADA ─────────────────────────────────────────────────
\AddToShipoutPicture*{%
    \put(54,700){\includegraphics[width=2cm]{\logoup}}
    \put(420,706){\includegraphics[width=5cm]{\logousach}}
}

\vspace*{0.5cm}
\begin{center}
    \color{azulFuerte}
    {\huge\textbf{GUÍA DE EJERCICIOS}}\\[0.5cm]
    {\Huge\textbf{\curso}}\\[0.3cm]
    {\Large\guiasnumero}\\[0.8cm]

    \begin{tcolorbox}[colback=celesteSuave, colframe=azulFuerte, arc=10pt,
        width=0.88\textwidth, center, boxrule=1.5pt]
        \centering
        \vspace{0.2cm}
        {\Large\textbf{Contenidos}}\\[0.4cm]
        \color{black}
        \begin{itemize}[leftmargin=3cm]
            \listacontenidos
        \end{itemize}
        \vspace{0.2cm}
        \smallskip
        {\small\color{azulFuerte}
        \nivelbasico\ Ej.\ 1–5 \quad
        \nivelintermedio\ Ej.\ 6–10 \quad
        {\color{rojoAlerta}\textbf{[Biyectividad]}} Ej.\ 11–13 \quad
        \nivelpep\ Ej.\ 14–15}
    \end{tcolorbox}
\end{center}

\vspace{1cm}
\begin{center}
    {\Large\textbf{\semestre}}\\[0.6cm]
    {\large\textbf{\preparadopor}}
\end{center}
\vspace{1.5cm}
\begin{center}
    {\color{azulFuerte}\hrule height 0.5pt}
    \vspace{0.3cm}
    \begin{minipage}{0.45\textwidth}
        \centering\textbf{Instagram:}\\
        \small\href{\instagraminstitucion}{@usach.premium}
    \end{minipage}
    \begin{minipage}{0.45\textwidth}
        \centering\textbf{Web:}\\
        \small\href{http://\webinstitucion}{\webinstitucion}
    \end{minipage}
\end{center}

\pagenumbering{arabic}
\newpage
\MarcaAgua

% ════════════════════════════════════════════════════════════
% SECCIÓN I — Dominio, Recorrido y Monotonía
% ════════════════════════════════════════════════════════════
\subsection*{I.\; Dominio, Recorrido y Crecimiento/Decrecimiento \nivelbasico}

\begin{enumerate}[label=\color{azulFuerte}\bfseries\arabic*., leftmargin=*, itemsep=1.2em]

% ── Ej 1 ─────────────────────────────────────────────────────
\item Determine el dominio natural de la función:
\[
    f(x) = \sqrt{x^{2}-9}
\]

% ── Ej 2 ─────────────────────────────────────────────────────
\item Determine el dominio natural y el recorrido de la función:
\[
    g(z) = \frac{1}{\sqrt{25-z^{2}}}
\]

% ── Ej 3 ─────────────────────────────────────────────────────
\item Determine el dominio natural de la función:
\[
    h(x) = \frac{\sqrt{x+4}}{x^{2}-1}
\]

% ── Ej 4 ─────────────────────────────────────────────────────
\item Determina los intervalos donde la función es creciente y donde es decreciente:
\[
    g(x) = 3(x-1)^{2}+2
\]

% ── Ej 5 ─────────────────────────────────────────────────────
\item Determina los intervalos donde la función es creciente y donde es decreciente:
\[
    h(x) = -(x+3)^{2}+5
\]

\end{enumerate}

\vspace{0.5cm}\newpage

% ════════════════════════════════════════════════════════════
% SECCIÓN II — Composición e Inyectividad
% ════════════════════════════════════════════════════════════
\subsection*{II.\; Composición e Inyectividad \nivelintermedio}

\begin{enumerate}[label=\color{azulFuerte}\bfseries\arabic*., leftmargin=*, itemsep=1.4em, start=6]

% ── Ej 6 ─────────────────────────────────────────────────────
\item Considere las funciones $f(x)=\sqrt{x+3}$ y $g(x)=2x-3$, ambas definidas en su dominio natural.
\begin{itemize}
    \item Encuentre el dominio de $(f\circ g)(x)$.
    \item Calcule $(f\circ g)(8)$.
    \item Determine la preimagen de $3$ vía $(f\circ g)$, si es que existe.
\end{itemize}

% ── Ej 7 ─────────────────────────────────────────────────────
\item Considere las funciones $f(x)=\sqrt{\dfrac{x+1}{x-2}}$ y $g(x)=x+3$, ambas en su dominio natural.
\begin{itemize}
    \item Determine el dominio de $(f\circ g)(x)$.
    \item Calcule $(f\circ g)(2)$.
    \item Determine la preimagen de $1$ vía $(f\circ g)$, si es que existe.
\end{itemize}

% ── Ej 8 ─────────────────────────────────────────────────────
\item Determine el valor de $k$ de modo que $f\circ g = g\circ f$, siendo:
\[
    f(x) = \frac{5x+3}{4} \qquad \text{y} \qquad g(x)=\frac{x+k}{5}
\]

% ── Ej 9 ─────────────────────────────────────────────────────
\item Considere la función $g:[2,+\infty[\,\to\mathbb{R}$ definida por
$g(x)=x^{2}-4x+8$.
\begin{itemize}
    \item Demuestre que $g$ es inyectiva en $[2,+\infty[$.
    \item Determine $\operatorname{Rec}(g)$ y argumente por qué $g$ \emph{no} es sobreyectiva.
\end{itemize}

% ── Ej 10 ────────────────────────────────────────────────────
\item Considere la función $f(x)=\sqrt{3-\sqrt{2x-1}}$.
\begin{itemize}
    \item Encuentre el dominio natural de $f$.
    \item Demuestre que $f$ es estrictamente decreciente en su dominio.
    \item Determine $\operatorname{Rec}(f)$.
\end{itemize}

\end{enumerate}

\vspace{0.5cm}\newpage

% ════════════════════════════════════════════════════════════
% SECCIÓN III — Biyectividad y Función Inversa
% ════════════════════════════════════════════════════════════
\subsection*{III.\; Biyectividad y Función Inversa \
    {\color{azulFuerte}\textbf{[Intermedio–Avanzado]}}}

\begin{enumerate}[label=\color{azulFuerte}\bfseries\arabic*., leftmargin=*, itemsep=1.4em, start=11]

% ── Ej 11 ────────────────────────────────────────────────────
\item Sea $f:[0,+\infty[\,\to B$ la función definida por:
\[
    f(x) = \frac{x}{x+2}
\]
\begin{itemize}
    \item Demuestre que $f$ es inyectiva.
    \item Determine el conjunto $B$ de modo que $f$ sea biyectiva.
    \item Encuentre $f^{-1}$ indicando su dominio y recorrido.
\end{itemize}

% ── Ej 12 ────────────────────────────────────────────────────
\item Sea $f:A\to B$ la función definida por $f(x)=-x^{2}+6x-5$, donde $3\in A$.
\begin{itemize}
    \item Determine los conjuntos $A$ y $B$ de modo que $f$ sea biyectiva.
    \item Encuentre $f^{-1}$ indicando su dominio y recorrido.
\end{itemize}

% ── Ej 13 ────────────────────────────────────────────────────
\item Sea $f:\,]2,+\infty[\,\to B$ la función definida por:
\[
    f(x) = \frac{x}{x-2}
\]
\begin{itemize}
    \item Determine $B$ de modo que $f$ sea biyectiva y demuéstrelo.
    \item Encuentre $f^{-1}$ indicando su dominio y recorrido.
\end{itemize}

\end{enumerate}

\vspace{0.5cm}\newpage

% ════════════════════════════════════════════════════════════
% SECCIÓN IV — Nivel PEP
% ════════════════════════════════════════════════════════════
\subsection*{IV.\; Ejercicios tipo PEP \nivelpep}

\begin{tcolorbox}[colback=rojoAlerta!8, colframe=rojoAlerta, arc=6pt,
    boxrule=0.8pt, left=4pt, right=4pt, top=3pt, bottom=3pt]
    \small\color{black}
    Los ejercicios siguientes corresponden al formato y dificultad de una PEP
    de Cálculo I. Se recomienda intentarlos sin ver el solucionario.
\end{tcolorbox}

\vspace{0.4cm}

\begin{enumerate}[label=\color{azulFuerte}\bfseries\arabic*., leftmargin=*, itemsep=1.6em, start=14]

% ── Ej 14 ────────────────────────────────────────────────────
\item Dadas las funciones $f(x)=\sqrt{x-3}$ y $g(x)=\sqrt{25-x^{2}}-1$,
    ambas definidas en su respectivo dominio natural:
\begin{itemize}
    \item Encuentre el dominio de $(f\circ g)(x)$.
    \item Halle explícitamente la regla de asignación de $(f\circ g)(x)$.
\end{itemize}

% ── Ej 15 ────────────────────────────────────────────────────
\item Considere la función
\[
    f(x) = \sqrt{\frac{4}{\sqrt{x+2}}-1}
\]
y la función lineal $g(x)=x+1$.
\begin{itemize}
    \item Halle el dominio natural de $f$.
    \item Demuestre que $f$ es estrictamente decreciente en su dominio.
    \item Determine $\operatorname{Rec}(f)$.
    \item Determine la función compuesta $f\circ g$, indicando expresamente
          su dominio y regla de asignación.
\end{itemize}

\end{enumerate}

\vspace{1cm}

% ════════════════════════════════════════════════════════════
% SOLUCIONARIO
% ════════════════════════════════════════════════════════════
\begin{soluciones}

\begin{solbox}{I. Soluciones — Dominio, Recorrido y Crecimiento/Decrecimiento}
\begin{enumerate}[label=\color{azulFuerte}\bfseries\arabic*., leftmargin=*, itemsep=0.6em]

    \item $\operatorname{Dom}(f) = (-\infty,-3]\cup[3,+\infty)$

    \item $\operatorname{Dom}(g)=(-5,5)$;\quad
          $\operatorname{Rec}(g)=\bigl[\tfrac{1}{5},+\infty\bigr)$

    \item $\operatorname{Dom}(h)=[-4,+\infty)\setminus\{-1,\,1\}$

    \item $g(x)=3(x-1)^2+2$: decreciente en $(-\infty,1]$, creciente en $[1,+\infty)$

    \item $h(x)=-(x+3)^2+5$: creciente en $(-\infty,-3]$, decreciente en $[-3,+\infty)$

\end{enumerate}
\end{solbox}

\vspace{0.4cm}

\begin{solbox}{II. Soluciones — Composición e Inyectividad}
\begin{enumerate}[label=\color{azulFuerte}\bfseries\arabic*., leftmargin=*, itemsep=0.6em, start=6]

    \item $\operatorname{Dom}(f\circ g)=[0,+\infty)$;\quad
          $(f\circ g)(8)=4$;\quad
          Preimagen de $3$: $x=\dfrac{9}{2}$

    \item $\operatorname{Dom}(f\circ g)=(-\infty,-4]\cup(-1,+\infty)$;\quad
          $(f\circ g)(2)=\sqrt{2}$;\quad
          Preimagen de $1$: no existe

    \item $k=-12$

    \item $g$ inyectiva (creciente en $[2,+\infty)$);\quad
          $\operatorname{Rec}(g)=[4,+\infty)\neq\mathbb{R}$ (no sobreyectiva)

    \item $\operatorname{Dom}(f)=\bigl[\tfrac{1}{2},5\bigr]$;\quad
          $f$ estrictamente decreciente;\quad
          $\operatorname{Rec}(f)=[0,\sqrt{3}]$

\end{enumerate}
\end{solbox}

\vspace{0.4cm}

\begin{solbox}{III. Soluciones — Biyectividad y Función Inversa}
\begin{enumerate}[label=\color{azulFuerte}\bfseries\arabic*., leftmargin=*, itemsep=0.7em, start=11]

    \item $B=[0,1)$;\quad $f$ biyectiva;\quad
          $f^{-1}(y)=\dfrac{2y}{1-y}$,\;
          $\operatorname{Dom}(f^{-1})=[0,1)$,\;
          $\operatorname{Rec}(f^{-1})=[0,+\infty)$

    \item $A=[3,+\infty)$,\; $B=(-\infty,4]$;\quad
          $f^{-1}(x)=3+\sqrt{4-x}$,\;
          $\operatorname{Dom}(f^{-1})=(-\infty,4]$,\;
          $\operatorname{Rec}(f^{-1})=[3,+\infty)$

    \item $B=(1,+\infty)$;\quad $f$ biyectiva;\quad
          $f^{-1}(y)=\dfrac{2y}{y-1}$,\;
          $\operatorname{Dom}(f^{-1})=(1,+\infty)$,\;
          $\operatorname{Rec}(f^{-1})=(2,+\infty)$

\end{enumerate}
\end{solbox}

\vspace{0.4cm}

\begin{solbox}{IV. Soluciones — Nivel PEP}
\begin{enumerate}[label=\color{azulFuerte}\bfseries\arabic*., leftmargin=*, itemsep=0.7em, start=14]

    \item $\operatorname{Dom}(f\circ g)=[-3,3]$;\quad
          $(f\circ g)(x)=\sqrt{\sqrt{25-x^{2}}-4}$

    \item $\operatorname{Dom}(f)=(-2,14]$;\quad
          $f$ estrictamente decreciente;\quad
          $\operatorname{Rec}(f)=[0,+\infty)$\\[4pt]
          $(f\circ g)(x)=\sqrt{\dfrac{4}{\sqrt{x+3}}-1}$,\quad
          $\operatorname{Dom}(f\circ g)=(-3,13]$

\end{enumerate}
\end{solbox}

\vspace{0.8cm}
\begin{center}
    \small\textbf{Usach Premium} — ¡Nos vemos en la ayudantía! \\
    \textit{``Por y para estudiantes.''}
\end{center}

\end{soluciones}

\end{document}
